\begin{table}[H]
\centering
\begin{tabular}{| c | c |}
\hline
Parameter(s) 											&	Asimov A 										\\
\hline
$\sin^{2}\theta_{23}$									&	0.528											\\
$\sin^{2} \theta_{13}$ reactors							&	0.0219										\\
$\sin^{2} \theta_{12}$									&	0.304											\\
$|\Delta m^{2}_{32}|$ (NH) / $|\Delta m^{2}_{31}|$ (IH)	&	$2.509 \times 10^{-3}~\mbox{eV}^2/\mbox{c}^4$	\\
$\Delta m^{2}_{21}$										&	$7.53 \times 10^{-5}~\mbox{eV}^2/\mbox{c}^4$	\\
$\delta_{CP}$											&	-1.601											\\
Mass Hierarchy											&	Normal											\\
$\beta$													&   1												\\  
\hline
\end{tabular}
\caption{
Values of oscillation parameters used to compute the event rates, systematic effects and sensitivity studies corresponding to the Asimov A data set, which is the MC expected distribution in a certain oscillation hypothesis.
In the Asimov data set A the nominal values of $\sin^2 2 \theta_{13}$, $\sin^2 2 \theta_{12}$
and $\Delta m^2_{21}$  are from \cite{pdg_2014}, while all the other oscillation parameter values corresponds to the most probable values obtained by the Bayesian analysis on the T2K run 1-4 neutrino mode data \cite{T2KLongPaper}. The mass hierarchy is not marginalised but fixed to either NH or IH.
}
\label{tab:osc_params_joint_asimovA}
\end{table}